\section{The Duty of the Wise Man}

Self-help books have a large presence in bookstores today, but there is nothing new about that. In the Ancient City, the way of life was fixed by ancestors and even when written down, they were known only to the few. However, as that system began to break up, there was no shortage in the ancient world of books telling men how to live. As circumstances changed, the books changed. In the medieval world, books like The Imitation of Christ were intended for the sacerdotal class. The aristocracy had many books on the proper code of conduct for chivalry.

When the Medieval structure collapsed, new types of books were published to accommodate the change in circumstances. This happened first in Renaissance Italy, and the events in Florence illustrate that. Following the failure of Girolamo Savonarola\footnote{\url{https://gornahoor.net/?p=2744}} to establish a theocracy in Florence, Niccolo Machiavelli took a more secular perspective, writing a self-help book for a Prince. As the moneyed patrician class squeezed out the landed aristocracy, there was the need for another type of self-help book. This was for the Courtier, who assisted at the court of the Prince. His qualities and behaviours were prescribed by \textbf{Baldesar Castiglione} in The Book of the Courtier. This book was immensely influential and was translated into several European languages. Its ripples can still be felt today.

Castiglione saw himself as recovering the virtues and character of the ancient Roman and Greek pagans. While, perhaps, there is little need for a Courtier today, it still has great value as an antidote to the crudities of popular culture today and as a guide to the development of a noble, or Aryan, character. Properly understood, moreover, it can be seen as a spiritual guide as the the pursuit of a life of virtue and the development of various skills can lead to greater self-awareness and self-mastery.

The courtier was expected to be skilled in the martial arts of weaponry and wrestling. He should be an expert horseman, hunter and know how to swim, jump, run and throw stones. But the courtier is not one-dimensional, as the arts are also important. He must be able to dance, to play a musical instrument, and to be able to draw or paint. This is a recovery of ancient virtues since both Plato and Aristotle expected a gentleman to play an instrument. The initiation ceremony of the \textbf{Kouretes} involved a dance while wielding weapons. The Courtier was expected to be quite literate and familiar with the great works. This would be reflected in his own style of writing and speaking.

The fundamental attitude of the Courtier was to do everything with a type of style and grace done with a certain nonchalance that made the difficult seem effortless. Castiglione coined a new word for this: \emph{Sprezzatura}. He explains:


\begin{quotex}
I have found quite a universal rule which in this matter seems to me valid above all others, and in all human affairs whether in word or deed: and that is to avoid affectation in every way possible as though it were some very rough and dangerous reef; and to practice in all things a certain sprezzatura, so as to conceal all art and make whatever is done or said appear to be without effort and almost without any thought about it.
\end{quotex}


This sprezzatura has spilled into popular culture as seen in, perhaps, \textbf{Humphrey Bogart} or the \textbf{James Bond} films. These characters are always accomplished and competent, but go about it in a ``cool'' way. They always utter the right phrase, that is at once apt, even self-deprecating. The era of Cassius Clay brought in an anti-sprezzatura of boastfulness and crudity, which equally persists today.

Castiglione exhorts us to live a unified life. Most of us are internally discordant, willing one thing today and its opposite tomorrow. To be a courtier requires a great deal of self-awareness. Today, the emphasis is on the \emph{elan vital}, where everyone ``lets it all hang out''. However, the courtier depends on the \emph{frein vital}, or the inner check in the sense of \textbf{Irving Babbitt}. This is the inner battle of personal order over personal chaos, where a man is driven to and fro by circumstances, whim, and instinct. If purity of heart is to will one thing, then the Courtier is pure in heart. This is how Castiglione describes it:

\begin{quotex}
Our Courtier must be cautious in his every action and see to it that prudence attends whatever he says or does. And let him take care not only that his separate parts and qualities be excellent, but that the tenor of his life be such that the whole may correspond to these parts, and may be seen to be, always and in everything, such as never to be discordant in itself, but form a single whole of all these good qualities; so that his every act may stem from and be composed of all the virtues which the Stoics hold to be the duty of the wise man.
\end{quotex}


Just to refresh the reader's memory, the cardinal Stoic virtues are: \textbf{wisdom, courage, justice, and temperance}. To achieve the virtuous life, the Stoic must bend his will to follow the cosmic order, or Logos. This requires a life of self-remembering. The disordered, unfree, and unwise mind wants a simple formula, so he can act the same way no matter what the particular circumstances may be. However, the Courtier must take everything into account because in the Whole, all the parts are interrelated. If a man takes the following advice of Castiglione seriously, he will develop virtue, nobility, self-awareness, a powerful will, and a sense of transcendence.

\begin{quotex}
Let him consider all what he does or says, the place where he does it, in whose presence, its timeliness, the reason for doing it, his own age, his profession, the end at which he aims, and the means by which he can reach it.
\end{quotex}

\flrightit{Posted on 2011-10-19 by Cologero}

